\documentclass[12pt]{article}
\usepackage[utf8]{inputenc}
\usepackage[T1]{fontenc}
\usepackage{lmodern}
\usepackage{microtype}
\usepackage[hidelinks]{hyperref}
\usepackage[a4paper, margin=3cm]{geometry}
\usepackage{parskip}

\title{Without a goal we cannot grow}
\date{09/06/26}

\begin{document}
\maketitle

Let's take a look at the grandiosity of Shri Mataji's vision and the current status of our Sahaja Yoga organizations.
On the one hand, She wants us to emancipate humanity, save all the humans, be perfect Sahaja Yogis, while on the other hand, we barely have any reach in society, nobody really knows of Sahaja Yoga, and we have no chance in a proper debate that involves actual problems of the world or society.
If one were a cynic, one could say we've failed.
Well, thank God we're not that.
So, let us look at why it is so difficult to achieve those goals.
First of all, a lot of work would be required to get there.
To grow spiritually needs serious effort and sacrifice - to change existing behaviours, purify the attention, and study the theory rigorously.
Second of all, there are a lot of other things that we have to do as well, like getting married, having children, taking care of the community.
Nevertheless, with the right focus and collective effort it should still be possible, or?
After all, this is something Shri Mataji wants from us.
So something that should be worth dying for, right?
If we look at the current operations in the Sahaja Yoga organizations, like the Stichting in NL, then it's even hard to identify that the vision of Shri Mataji we mentioned earlier plays a central role.
In reality, the problem is that our goal is not actually to grow spiritually, that is, to achieve Nirvikalpa.
Our goal is also not to emancipate humanity.
At the moment we're rather trying to keep Sahaja Yoga alive and keep its activities running.
In my opinion that is the problem.
Without setting the right goal it is no surprise that we're not reaching the target.
We have to understand that the motivation for doing things - especially difficult things - comes achieving a goal.
E.g. if I want to clean my house I'll find myself having difficulty to find the time or to prioritze it because it is in itself not a very fruitful activity.
It is something we do because it's good to do but not because it's particularly rewarding.
That is also why we call it a chore.
However, if we have guests over suddenly we find ourselves cleaning without any argument or debate whether we feel like doing it or not.
Suddenly, we see it as a step towards our higher goal and we just work through it like ticking off a box on our checklist.
The same goes for many other things.
I believe we can all think of tasks that we would really like to do but we just don't find the time, or space to do them.
All those tasks are left hanging because there is no project that we need it for or because it is not part of any goal that we have set for ourselves.
Our individual meditation is one of those things.
We all have best intentions to meditate deeply and sincerly.
However, when life gets busy we suddenly drop it, or we don't find the time anymore to go deeper.
Or when we fail to meditate properly we also don't face any consequences.
We can still join the collective meditation, or participate in Pujas.
There is no accountability, at least not on a collective level.
It doesn't matter at all to the collective goals if we don't meditate sometimes, or if our growth is slow for some time.
That should be different in my opinion.
Image we all set ourselves the goal to reach Nirvikalpa by the end of the year.
Suddenly, we'd have to spend a significant portion of our time for our meditation because we'd have to become perfect in our meditation.
We'll be studying Shri Mataji's speeches, doing more Sahaj work and ultimately motivate each other to reach there.
Because after all we set that goal all together.
Or take the emancipation of humanity.
If our goal was really to emancipate humanity we wouldn't be doing Sahaja Yoga so half-heartedly.
We would feel compelled to understand the shortcomings of our society and how to solve them.
We'd be studying all the religions of the world, philosophy, science, and politics.
We'd be engaging with people from all walks of life to understand how we can help them.
Even if these are just far fetched examples.
One thing is for sure.
We surely wouldn't have time to waste arguing with each other, judging each other, fighting over petty little things, or hiding from the outside world in our beautiful Sahaj bubble, where we can just neglect all problem because "Shri Mataji is going to take care of everything."



% --- EDITORIAL FEEDBACK (10/06/26) ---
%
% WHAT WORKS
% - Core argument is clear and real: you've confused maintenance for mission.
% - Cleaning house metaphor lands — concrete, followable, illustrates the mechanism well.
% - The ending is the most alive writing in the piece. The whole piece should build toward it.
%
% WHAT DOESN'T WORK
% - No thesis up front. The real argument arrives at line 28. Everything before it is
%   throat-clearing. Cut or compress to two sentences, then open with the argument.
% - Prose is too loose. "or?", "right?", "like the Stichting in NL" — conversational
%   hedges that read as unconfident in writing. Remove them.
% - You undercut your own examples. "Even if these are just far fetched examples." Either
%   commit to them or find better ones. The hesitation kills the momentum.
% - The ending diagnoses but doesn't direct. The reader finishes without knowing what
%   you're actually asking for. Set a collective goal? Restructure accountability?
%   Strong diagnosis needs a concrete prescription.
% - Typo: line 50 "Image" should be "Imagine."
%
% THREE PRIORITIES
% 1. Open with a thesis: e.g. "Our problem is not capability or commitment.
%    It is that we have replaced mission with maintenance."
% 2. Cut everything before the argument begins.
% 3. End with a concrete ask, not just an accusation.
%
% The bones are good. Needs compression and a sharper spine.

\end{document}
